% Generated draft: requires mathematical and source review.
\documentclass[11pt]{article}
\usepackage[margin=25mm]{geometry}
\usepackage{fontspec}
\setmainfont{Latin Modern Roman}
\usepackage{amsmath,amssymb,mathtools,unicode-math}
\setmathfont{Latin Modern Math}
\usepackage{enumitem,needspace,xurl,hyperref}
\setlength{\parindent}{0pt}
\setlength{\parskip}{5pt}
\setlength{\emergencystretch}{3em}
\newcommand{\N}{\mathbb{N}}
\newcommand{\Z}{\mathbb{Z}}
\newcommand{\Q}{\mathbb{Q}}
\newcommand{\R}{\mathbb{R}}
\newcommand{\C}{\mathbb{C}}
\newcommand{\F}{\mathbb{F}}
\newcommand{\A}{\mathbb{A}}
\newcommand{\PP}{\mathbb{P}}
\newcommand{\E}{\mathbb{E}}
\newcommand{\eps}{\varepsilon}
\newcommand{\dd}{\,\mathrm{d}}
\begin{document}
\section*{UnsolvedMath ID: 10400215}
\textbf{Source problem code:} catalogue entry
\subsection{Definitions and mathematical statement}
Write \(s(M)\) for the shadow number of a \(3\)-manifold \(M\), as intended in the source. Its precise definition and the admissible shadows are not included in the supplied excerpt. Write \(\|M\|\) for the Gromov norm. For a closed oriented manifold, the usual simplicial-volume convention is
\[
 \|M\|=\inf\left\{\sum_i|a_i|:
 \sum_i a_i\sigma_i\text{ is a real singular cycle representing }[M]\right\}.
\]
A relative or other convention would need to be specified for manifolds outside that class.

\emph{Conjecture (D. Thurston).} The shadow number is quasi-linear in the Gromov norm, in the explicit multiplicative sense stated in the source: there exist constants \(c_1,c_2\), independent of \(M\), such that
\[
 c_1\|M\|\leq s(M)\leq c_2\|M\|
 \qquad\text{for every }3\text{-manifold }M
\]
in the intended domain. The natural nonvacuous reading of the comparison requires \(0<c_1\leq c_2<\infty\); the excerpt itself says only ``constants'' and does not state their signs. It also does not fix compactness, connectedness, orientation, or boundary hypotheses for ``any \(3\)-manifold.''

The displayed inequalities are preserved exactly as multiplicative bounds, with no additive error. In particular, in any domain where \(s(M)\geq0\), the upper bound forces \(s(M)=0\) whenever \(\|M\|=0\). This zero-norm consequence is part of the supplied formula and must be checked against the intended shadow-number definition; a restriction to hyperbolic manifolds or an additive correction is not silently introduced.
\subsection{Short English statement}
Thurston conjectures universal multiplicative upper and lower bounds comparing shadow number with Gromov norm. The source formula has no additive term and includes every manifold in its unspecified domain, so zero Gromov norm would force zero shadow number.
\subsection{Sources}
\begin{itemize}
\item Local UnsolvedMath registry, permanent ID 10400215, source code catalogue entry; attribution D. Thurston. Supplied title is truncated after ``Conjecture 12.10 --- (D.''.
\item Ohtsuki, \emph{Problems on invariants of knots and 3-manifolds} (2002), Conjecture 12.10, PDF page 164. Supplied URL: \url{https://msp.org/gtm/2002/04/gtm-2002-04-024s.pdf}; recorded access 2026-07-29; extraction pdf-font-structure. The collection distinguishes 231 numbered bold problem, conjecture, and question headings.
\item The registry labels the item open and OPEN-TRIAGE and describes the comparison as quasi-linear. Background status evidence remains NEEDS\_REVIEW pending a current primary-literature audit, and redistribution terms for the publicly accessible author/publisher copy require review.
\end{itemize}
\end{document}
